(1) 一维连续Dirac哈密顿量H=-i hbar v sigma_x partial_x+Delta(x) sigma_z，v>0，Delta(x)=Delta0 tanh(x/ell)，Delta0,ell>0。证明存在与sigma_y相关的手征反对易对称性。
(2) 求零能本征态的两分量微分方程及可归一化局域解，计算其空间包络、无穷远局域长度。
(3) 平方哈密顿量获得两类手征通道对应的Schrödinger型伴随势，证明零模所在通道的有效势可含负的畴壁束缚阱，而H²整体非负。
(4) 把质量畴壁方向反转或让Delta不改变符号，分析零模手征翻转或消失的条件；讨论破坏手征对称的常量单位阵势与边界效应。
